% ---
% school: "college"
% level: "8eme"
% section: null
% chapter: "العمليات والحساب على الأعداد الكسرية"
% source: "Axiomaths"
% title: "تفكيك ومقارنة"
% country: "Tunisie"
% year: 2026
% professor: "Axiomaths"
% difficulty: "Moyen"
% points: null
% tags: ["fractions", "factorisation", "comparaison", "expressions-litterales"]
% ---
% Exercice 10/18 de la série 8ème « العمليات والحساب على الأعداد الكسرية » (statement + \axsol + solution). Compiler avec XeLaTeX + axiomath.sty (+ axiomaths-logo.png).
\documentclass[12pt,a4paper]{article}
\usepackage[level=8b,color=15803D]{axiomath}
\axsetup{title={العمليات والحساب على الأعداد الكسرية}, corrige}

\begin{document}

\begin{corr}[title={تفكيك ومقارنة}, stars=3, source={Axiomaths}]
$a$ و $b$ عددان كسريان. نعتبر $P=a\times b+a$ و $Q=a\times b+b$.
\begin{enumerate}
  \item فكّك $P$ و $Q$ إلى جذاءين.
  \item احسب $P$ و $Q$ من أجل $a=\dfrac{4}{3}$ و $b=\dfrac{1}{3}$ ثم قارنهما.
  \item بين أن $P-Q=a-b$.
  \item استنتج أنه إذا كان $a>b$ فإن $P>Q$ ، وإذا كان $a<b$ فإن $P<Q$.
  \item قارن $P$ و $Q$ من أجل $a=-\dfrac{2}{3}$ و $b=-\dfrac{1}{2}$ دون حساب $P$ و $Q$ ، ثم تحقق بالحساب.
\end{enumerate}
\axsol
\begin{enumerate}[itemsep=14pt]
  \item\leavevmode
  \begin{prop}[التفكيك إلى جذاء]
  نبحث عن العامل المشترك : \ib{$k\times x+k\times y=k\times(x+y)$}.
  \end{prop}
  في $P=a\times b+a\times1$ العامل المشترك هو $a$ ، وفي $Q=a\times b+b\times1$ العامل المشترك هو $b$ :
  \[P=a\times(b+1)\qquad\qquad Q=b\times(a+1)\]
  \item نعوّض :
  \[P=\dfrac{4}{3}\times\left(\dfrac{1}{3}+1\right)=\dfrac{4}{3}\times\dfrac{4}{3}=\dfrac{16}{9}\]
  \[Q=\dfrac{1}{3}\times\left(\dfrac{4}{3}+1\right)=\dfrac{1}{3}\times\dfrac{7}{3}=\dfrac{7}{9}\]
  للعددين نفس المقام $9$ و $16>7$ ، إذن \ib{$P>Q$}.
  \item نطرح، وعند حذف الأقواس المسبوقة بعلامة $(-)$ تتغير علامات ما بداخلها :
  \[P-Q=(a\times b+a)-(a\times b+b)=a\times b+a-a\times b-b\]
  الحدان $a\times b$ و $-a\times b$ يختزلان :
  \[\ib{$P-Q=a-b$}\]
  \begin{prop}[المقارنة بالفرق]
  \ib{$x>y$} يعني أن \ib{$x-y>0$}.
  \end{prop}
  \item $P-Q=a-b$. إذا كان $a>b$ فإن $a-b>0$ ، إذن $P-Q>0$ أي $P>Q$.
  وإذا كان $a<b$ فإن $a-b<0$ ، إذن $P-Q<0$ أي $P<Q$.
  \item نقارن $a$ و $b$ أولا بتوحيد المقامين :
  \[a=-\dfrac{2}{3}=-\dfrac{4}{6}\qquad b=-\dfrac{1}{2}=-\dfrac{3}{6}\]
  $-4<-3$ ، إذن $a<b$ ، ومن السؤال 4 : \ib{$P<Q$}.
  التحقق بالحساب :
  \[P=\left(-\dfrac{2}{3}\right)\times\left(-\dfrac{1}{2}+1\right)=-\dfrac{2}{3}\times\dfrac{1}{2}=-\dfrac{1}{3}=-\dfrac{2}{6}\]
  \[Q=\left(-\dfrac{1}{2}\right)\times\left(-\dfrac{2}{3}+1\right)=-\dfrac{1}{2}\times\dfrac{1}{3}=-\dfrac{1}{6}\]
  و $-2<-1$ ، فعلا $P<Q$.
\end{enumerate}
\textbf{الفكرة :} لمقارنة عبارتين نطرح إحداهما من الأخرى : إشارة الفرق تكفي، ولا حاجة لحساب كل منهما.
\end{corr}

\end{document}
